\documentclass[11pt]{article}
\usepackage[margin=1.1in]{geometry}
\usepackage{amsmath,amssymb,amsthm}
\usepackage{hyperref}

\newtheorem{theorem}{Theorem}[section]
\theoremstyle{definition}
\newtheorem{definition}[theorem]{Definition}
\newtheorem{conjecture}[theorem]{Conjecture}
\theoremstyle{remark}
\newtheorem*{remark}{Remark}

\title{A Machine-Verified Disproof of TLMC Conjecture \#1556:\\ The Chromatic Number of $G(\mathbb{Z}^2, \{1,2,4\})$}
\author{Batch 9 solution submission}
\date{September 12, 2026}

\begin{document}
\maketitle

\begin{abstract}
We disprove a conjecture of the Last Math Competition asserting that the chromatic number of the planar distance graph $G(\mathbb{Z}^2, D)$ with $D = \{1, 2, 4\}$ equals $7$. In fact $\chi = 3$: the map $c(x,y) = (x+y) \bmod 3$ is a proper $3$-coloring, and the axis triangle $\{(0,0),(1,0),(2,0)\}$ is a clique. The coloring is verified in Lean~4 by a mod-$3$ argument: lattice vectors of squared Euclidean length $1$, $4$ or $16$ have coordinate sum not divisible by $3$, since squares are $0$ or $1 \pmod 3$ while $2x^2 \not\equiv 1 \pmod 3$.
\end{abstract}

\section{The conjecture}

\begin{conjecture}[TLMC-00000001556]\label{conj:1556}
The chromatic number of the planar distance graph $G(\mathbb{Z}^2, D)$ with $D = \{1, 2, 4\}$ is $7$.
\end{conjecture}

\section{The disproof}

\begin{theorem}\label{thm:1556}
Conjecture~\ref{conj:1556} is \textbf{false}: $\chi(G(\mathbb{Z}^2, \{1,2,4\})) = 3$.
\end{theorem}

\begin{proof}
\emph{Upper bound.} Define $c(x, y) = (x + y) \bmod 3$. Two lattice points at Euclidean distance $d \in \{1,2,4\}$ differ by a vector $(a,b)$ with $a^2 + b^2 = d^2 \in \{1, 4, 16\}$. Suppose $3 \mid (a+b)$; then $b \equiv -a \pmod 3$, so
\[
a^2 + b^2 \equiv 2a^2 \pmod 3,
\]
and $2a^2 \in \{0, 2\} \pmod 3$ since squares are $0$ or $1 \pmod 3$. But $a^2 + b^2 \equiv d^2 \equiv 1 \pmod 3$ (as $1, 4, 16 \equiv 1$), a contradiction. Hence the color changes along every edge, and $c$ is a proper $3$-coloring.

\emph{Lower bound.} The points $(0,0)$, $(1,0)$, $(2,0)$ are pairwise at distances $1, 1, 2 \in D$, forming a triangle; hence $\chi \ge 3$.

Together $\chi = 3 \ne 7$.
\end{proof}

\begin{remark}
The mistake in the conjecture is presumably that the relevant ``hard'' instances of lattice distance graphs involve non-axis distances (e.g.\ $D = \{1, \sqrt{2}, 2\}$ or $\{1, 2, 5\}$-type sets with knight moves); for the purely axis-parallel set $\{1,2,4\}$ every edge is axis-parallel and the mod-$3$ diagonal coloring applies.
\end{remark}

\section{Machine verification}

\texttt{Batch9.lean} (namespace \texttt{TLMCBatch9}) contains \texttt{two\_sq\_ne\_one} ($2x^2 \ne 1$ in $\mathbb{Z}/3$), the arithmetic core \texttt{key\_arith}, the coloring \texttt{col} with correctness \texttt{col\_proper}, the clique \texttt{triangle\_clique}, and the packaged refutation \texttt{tlm1556} with \texttt{tlm1556\_false}. Compilation via \texttt{lake env lean Batch9.lean} exits $0$ with zero errors and warnings; \texttt{\#print axioms} reports the standard three axioms at most (\texttt{tlm1556\_false} depends on none); no \texttt{sorry} appears.

\end{document}
